% Chapter VII, Section 74 and Bibliographical Notes
% Chandrasekhar, "Hydrodynamic and Hydromagnetic Stability"
% Transcribed from scanned PDF

%% ============================================================
%% SECTION 74
%% ============================================================
\section{Experiments on the stability of viscous flow between rotating cylinders}\label{sec:7-74}

The first experiments on the onset of instability in the viscous flow
between rotating coaxial cylinders are those of Taylor.  Experiments
similar to Taylor's, and others related to the same phenomenon, have
since been carried out by Lewis, Terada and Hattori, Wendt, Taylor
himself, Schultz-Grunow and Hein, Donnelly, and Donnelly and Fultz.
We shall restrict our account of the experiments to those of Donnelly,
and Donnelly and Fultz as they are not only typical of the class, but they
are also among the most precise and the most extensive; moreover, they
bear directly on certain aspects of the theory considered in this chapter.

First, we may observe that the experiments one can perform in the
present connexion are of two kinds: those which involve the measure\-ment
of the torque exerted on the outer cylinder (kept at rest) for different
rates of rotation of the inner cylinder, and those in which the motions in
the fluid are observed by suitable tracers.

In the experiments of the first kind, the onset of instability is detected
by a principle not unlike the Schmidt--Milverton principle for the
detection of the onset of thermal instability (cf.\ \S\,18\,($b$)).  In the present
instance, the torque exerted on the outer cylinder (which is suitably
suspended) is measured as a function of the angular velocity $\Omega_1$ of the
inner cylinder.  It is known that so long as the flow between the cylinders
is laminar and conforms to the law (1), the relation is a linear one and is
given by
\begin{equation}\label{eq:7-321}
\text{torque} = 4\pi\mu\,\frac{R_1^2\,R_2^2\,H}{R_2^2 - R_1^2}\,\Omega_1,
\end{equation}
where $H$ is the height of the suspended cylinder.  At the onset of instability,
the effective viscosity will abruptly start increasing, and this
will be reflected in a sudden break in the relation~\eqref{eq:7-321}.

In the experiments of the second kind, the onset of instability is
detected by visual observations of sudden changes in the character of the
fluid motion occurring as the relative angular velocity of the two
cylinders is increased.  While these experiments are not as precise as
those on the torque measurements for determining the critical angular
velocities at which instability arises, they are not restricted to any special
value of $\mu$ ($= \Omega_2/\Omega_1$) such as $\mu = 0$.  Indeed, the important range of
negative $\mu$ is readily accessible to observations only by this second
method.

%% ------------------------------------------------------------
%% (a) Critical Taylor numbers by torque measurements
%% ------------------------------------------------------------
\subsection{The determination of the critical Taylor numbers, for the case $\mu = 0$,
by torque measurements}\label{sec:7-74a}

The most precise experiments on the determination of the critical
Taylor numbers for the onset of instability, for the case $\mu = 0$, are those
of Donnelly.

In his experiments, Donnelly used a rotating cylinder viscometer (see
Fig.~77) in which only a central section of the outer cylinder was sus\-pended;
the fixed end sections served as a guard cylinder for eliminating
end effects.  The viscometer was so designed that if the suspended
cylinder departed by so much as 0$\cdot$005~in., in any direction from exact
concentricity with the guard cylinder, the suspension would not swing
freely.  The system was thus protected from external vibrations and
lateral motions.

\figurePlaceholder{77}{Sketch of the viscometer used in Donnelly's experiments.}

The suspension system for the outer cylinders consisted of the follow\-ing
parts.  The straps by which the outer cylinder was carried were
attached to a long quartz rod (0$\cdot$018~in.\ in diameter and 11~cm long) at
the end of which was a torsion fibre (0$\cdot$005~in.\ in diameter and 10~cm long).
The fibre was in turn attached to a shaft whose angular position could be
read on a dial provided with a vernier.  A small mirror affixed to the
torsion fibre was directly in front of a stationary mirror.  The torque
exerted on the suspended cylinder was measured by a null method by
turning the dial so as to bring into coincidence the image of a cross-hair
(observed through an auto-collimator) reflected in both mirrors.  The
angular deflexion necessary to bring about the coincidence is a measure
of the torque.

The inner cylinder was driven by a shaft coupled to a motor and gear
reducing system.  The period of rotation was determined photo\-electrically
by using suitable scaling circuits.  In the experiments, the
speed of rotation was maintained constant to well within 0$\cdot$1 per cent.

The viscometer provided two interchangeable inner cylinders: one of
them left a gap width ($d/R_2 \sim 0{\cdot}05$) sufficiently small for the theory
of \S\,71 to be applicable; and the other left a gap width equal to the radius
of the inner cylinder so that the results of the experiments could be
compared with the calculations of \S\,73\,($b$).  The relevant dimensions
of the viscometer are given in Table~XXXVI.

To maintain the temperature of the liquid at a constant value during
the experiments, the viscometer was placed in a beaker kept in a constant
temperature bath.

The liquids used in the experiments and their relevant properties are
listed in Table~XXXVII.

\begin{table}[ht]
\centering
\caption{The dimensions and constants of Donnelly's viscometer}
\label{tab:XXXVI}
\medskip
\begin{tabular}{ll}
\hline
\multicolumn{2}{c}{\textit{Narrow gap}} \\
\hline
$R_1$ & $= 1{\cdot}89935 \pm 0{\cdot}0001$~cm \\
$R_2$ & $= 2{\cdot}00023 \pm 0{\cdot}0001$~cm \\
$H$   & $= 4{\cdot}9987 \pm 0{\cdot}0003$~cm \\
$C$   & $= 3{\cdot}93 \times 10^{-4}$ \\
\hline
\multicolumn{2}{c}{\textit{Wide gap}} \\
\hline
$R_1$ & $= 0{\cdot}99963 \pm 0{\cdot}0001$~cm \\
$R_2$ & $= 2{\cdot}00023 \pm 0{\cdot}0001$~cm \\
$H$   & $= 4{\cdot}9987 \pm 0{\cdot}0003$~cm \\
$C$   & $= 0{\cdot}967 \times 10^{-4}$ \\
\hline
\end{tabular}
\end{table}

\begin{table}[ht]
\centering
\caption{List of fluids and their physical properties}
\label{tab:XXXVII}
\medskip
\begin{tabular}{lccc}
\hline
Fluid & Temperature (${}^{\circ}$C) & $\rho$ (g) & $\nu$ (cm$^2$/sec) \\
\hline
CCl$_4$                          & 25$\cdot$00 & 1$\cdot$585 & $(1{\cdot}756 \pm 0{\cdot}03) \times 10^{-2}$ \\
CCl$_4$                          & 20$\cdot$95 & 1$\cdot$591 & $(6{\cdot}009 \pm 0{\cdot}03) \times 10^{-3}$ \\
CS$_2$                           & 25$\cdot$00 & 1$\cdot$255 & $(1{\cdot}935 \pm 0{\cdot}02) \times 10^{-2}$ \\
oil I, lot 111$^{\dag}$          & 25$\cdot$00 & 0$\cdot$8244 & $1{\cdot}330 \times 10^{-1}$ \\
oil D, lot 111$^{\dag}$          & 25$\cdot$00 & 0$\cdot$2813 & $2{\cdot}347 \times 10^{-2}$ \\
\hline
\end{tabular}

\smallskip
{\footnotesize $^{\dag}$ These oils are provided by the National Bureau of Standards (Washington, D.C.)
for calibrating viscometers. The National Bureau of Standards does not state the limits
of error on their measurements.}
\end{table}

If $K$ is the torsion constant of the fibre and $\Phi$ is the deflexion in radians,
then according to equation~\eqref{eq:7-321},
\begin{equation}\label{eq:7-322}
K\Phi = 4\pi\mu\,\frac{R_1^2\,R_2^2\,H}{R_2^2 - R_1^2}\,\Omega_1.
\end{equation}
In Donnelly's apparatus, the dial attached to the torsion fibre was
divided into 1,000 parts so that if $\phi$ is the measured deflexion in dial
divisions,
\begin{equation}\label{eq:7-323}
\nu = \left(\frac{K}{4\pi}\,\frac{R_2^2 - R_1^2}{R_1^2\,R_2^2\,H} \times 10^{-3}\right)\phi P
= C\,\phi\,P \quad (\text{say}),
\end{equation}
where $P$ ($= 2\pi/\Omega_1$) is the period of rotation of the inner cylinder.
Equation~\eqref{eq:7-323} applies when the flow is laminar; when it is not, we may
still regard the quantity on the right-hand side as giving an \textit{effective
viscosity}.

The value of $C$ was determined by calibration with liquids of known
viscosity rather than by an absolute determination of the torsion
constant~$K$.  The values of $C$ given in Table~XXXVI were obtained in
this manner.

In carrying out the experiments, the speed of rotation of the inner
cylinder was slowly increased from the lowest setting, keeping the
outer cylinder near its null position all the time.

\medskip
\noindent (i) \textit{Results of the experiments with the narrow gap}

The definition of the Taylor number appropriate for this case is (see
equation~(198))
\begin{equation}\label{eq:7-324}
T = \frac{4A\Omega_1}{\nu^2}\,d^4,
\end{equation}
where $A$ is given in equation~(145).  For $\mu = 0$, the explicit form of $T$ is
\begin{equation}\label{eq:7-325}
T = \frac{4\,R_1^2\,d^4}{\bigl(R_2^2 - R_1^2\bigr)}\left(\frac{\Omega_1}{\nu}\right)^{\!2}.
\end{equation}
According to the results given in Table~XXXIII, the critical Taylor
number for $\mu = 0$ is $3{\cdot}390 \times 10^3$.  The corresponding critical value of $\Omega_1$
is given by
\begin{equation}\label{eq:7-326}
\Omega_1(\text{critical}) = \left(3{\cdot}390 \times 10^3\,\frac{R_2^2 - R_1^2}{4\,R_1^2\,d^4}\right)^{1/2}\!\nu.
\end{equation}
For the values of $R_1$ and $R_2$ given in Table~XXXVI,
\begin{equation}\label{eq:7-327}
\Omega_1(\text{critical}) = 9{\cdot}448 \times 10^2\,\nu.
\end{equation}
For carbon tetrachloride at $25^{\circ}$~C instability is predicted for a period of
rotation
\begin{equation}\label{eq:7-328}
P(\text{critical}) = 1{\cdot}147 \pm 0{\cdot}006~\text{s (calculated)},
\end{equation}
where the limits of error are due to the uncertainty in $\nu$.

The results of Donnelly's experiments are shown in Fig.~78\,$a$, $b$.  In
these figures, $\phi P$ is plotted as a function of $P$.  We observe that as the
period of rotation of the inner cylinder is decreased, the viscosity as
given by equation~\eqref{eq:7-323} begins to increase abruptly at a certain period.
The break occurs with extreme sharpness; how sharply, can be seen from
Fig.~78\,$b$ in which the measurements in the immediate vicinity of the
critical period are exhibited.

The critical period of rotation determined, as intermediate between
the last of the stable and the first of the unstable points of observation, is
\begin{equation}\label{eq:7-329}
P(\text{critical}) = 1{\cdot}1337 \pm 0{\cdot}00095~\text{s (measured)}.
\end{equation}
This agrees very well with the calculated value~\eqref{eq:7-328}.

\figurePlaceholder{78a}{Plot of $\phi P$ (which is proportional to the effective viscosity)
as a function of the period of rotation, $P$, of the inner cylinder for the
narrow-gap case: $R_1 = 1{\cdot}9$~cm, $R_2 = 2{\cdot}0$~cm.  The liquid was carbon
tetrachloride at $25^{\circ}$~C\@.  $C$ is the calculated value for instability.}

\figurePlaceholder{78b}{Detail of measurements, shown in Fig.~78\,$a$,
about the critical speed.  The onset of instability is
characterized by a discontinuity both in the magnitude
of $\phi P$ and in the slope of the effective viscosity curve.}

\medskip
\noindent (ii) \textit{Results of the experiments with the wide gap} ($\eta = \frac{1}{2}$)

The definition of $T$ appropriate for this case is (cf.\ equation~(314))
\begin{equation}\label{eq:7-330}
T = \frac{64}{9}\!\left(\frac{\Omega_1\,R_1^2}{\nu}\right)^{\!2}.
\end{equation}
According to the results given in Table~XXXV, the critical Taylor
number for this case ($\kappa = 1$ and $\mu = 0$) is $3{\cdot}310 \times 10^4$.  The corresponding
critical value of $\Omega_1$ is given by
\begin{equation}\label{eq:7-331}
\Omega_1(\text{critical}) = \left(\frac{9 \times 3{\cdot}310 \times 10^4}{64\,R_1^4}\right)^{1/2}\!\nu.
\end{equation}
For the value of $R_1$ given in Table~XXXVI,
\begin{equation}\label{eq:7-332}
\Omega_1(\text{critical}) = 68{\cdot}28\,\nu.
\end{equation}
The experiments with this wide-gap viscometer were carried out with
an oil (oil~I in Table~XXXVII) supplied by the National Bureau of
Standards for calibrating viscometers.  For this oil, instability of flow
is predicted for a period of rotation
\begin{equation}\label{eq:7-333}
P(\text{critical}) = 0{\cdot}7508 \pm 0{\cdot}0004~\text{s (calculated)}.
\end{equation}
The results of Donnelly's experiments in this case are shown in Fig.~79\,$a$, $b$.  We observe that, again, an increase in the viscosity as given by
equation~\eqref{eq:7-323} begins abruptly at a certain critical period of rotation of

\figurePlaceholder{79a}{Plot of $\phi P$ as a function of $P$ for the wide-gap
case: $R_1 = 1{\cdot}0$~cm, $R_2 = 2{\cdot}0$~cm.  The liquid was oil~I at
$25^{\circ}$~C\@.  $C$ is the calculated value for instability.}

\figurePlaceholder{79b}{Detail of measurements, shown in Fig.~79\,$a$, near
the critical speed.  The onset of instability is characterized
by a discontinuity in the curve, but less pronounced than in the
narrow-gap case (Fig.~78\,$b$).  Curves 1 and 2 were taken
on different days.  $C$ is the calculated value for instability.}

\noindent the inner cylinder.  The measurements in the immediate vicinity of the
critical period exhibited in Fig.~79\,$b$ show how sharply it begins.

The two curves (1) and (2) in Fig.~79\,$b$ show that, while the magnitude
of $\phi P$ at the critical period varies from one sequence of measurements to
another, the critical period itself is reproducible; the value of the critical
period as experimentally determined is
\begin{equation}\label{eq:7-334}
P(\text{critical}) = 0{\cdot}7543 \pm 0{\cdot}0003~\text{s (measured)}.
\end{equation}
Again, this agrees very well with the calculated value~\eqref{eq:7-333}.

%% ------------------------------------------------------------
%% (b) Visual and photographic observations
%% ------------------------------------------------------------
\subsection{The dependence of the critical Taylor number on $\Omega_2/\Omega_1$. The results of
visual and photographic observations}\label{sec:7-74b}

While the experiments described in \S\,(a) above verify the theoretical
predictions regarding the critical Taylor numbers for the onset of
instability for the particular case $\mu = 0$, they do not provide
confirmation for other important aspects of the theory such as the
dependence on $\mu$ ($= \Omega_2/\Omega_1$) of the critical Taylor number and of the
wave number of the associated disturbance.  For these latter purposes,
it is essential that experiments are carried out in an apparatus in which
the two cylinders can be rotated independently of each other.
The first experiments in this category are those of Taylor.  Taylor's
experiments were carried out with an apparatus with a narrow enough
gap for the results of \S\,71\,($b$) to be applicable.  They have been repeated
by Donnelly and Fultz with an apparatus in which the gap was equal to
the radius of the inner cylinder so that the results of \S\,73\,($b$) are
similarly applicable.

Before presenting the results of Taylor's and Donnelly and Fultz's
experiments, we shall give a brief description of the apparatus used by
the latter authors (see Fig.~80).  The outer cylinder was a precision
bore Pyrex tube 94~cm long and inside diameter $R_2 = 6{\cdot}2846 \pm 0{\cdot}0006$~cm.
The inner cylinder was a brass tube turned to a radius
$R_1 = 3{\cdot}1432 \pm 0{\cdot}0001$~cm.  The Pyrex cylinder rested on a groove cut
in a plate which fitted on a brass base mounted on a rotating table.  The
brass cylinder was connected by means of a drive pin and locking collar
to a shaft which could be rotated independently of the brass base.  The
cylinders were rotated by separate pulleys and by two motors with
variable speed transmission.  The two cylinders were aligned to be coaxial
by suitable precision ball bearings.

\figurePlaceholder{80}{Sketch of the apparatus used in the experiments of Donnelly and
Fultz.  The ratio of the radii of the two cylinders is one-half.}

The motion of the fluid was traced by ink emitted from three small
holes at the middle of the inner cylinder.  The ink was a solution of
Nigrosine dye in the liquid used in the experiments and was of such
strength that drops of it neither sank nor rose in the liquid under the
experimental conditions.

The liquids used in the experiments were distilled water and glycerine
water mixtures of different strengths.  Kinematic viscosities in the range
$0{\cdot}01$--$0{\cdot}20$~cm$^2$/sec were thus available for the experiments.

In the experiments, the angular velocity of the outer cylinder was
treated as an independent variable.  The outer cylinder was first brought
to a certain speed and was, thereafter, kept at that speed.  After the
conditions had become stationary, the inner cylinder was set in rotation
and its speed was gradually increased until it was within a few per cent of
the expected critical value (as determined by preliminary observations).
A small amount of ink was then injected and allowed to spread over a
small region of the inner cylinder.  The speed of the inner cylinder was
then increased, slowly and carefully, until a definite sustained three-dimensional
motion was observed which indicated the onset of instability.

%% ------------------------------------------------------------
%% (i) Wave number observations
%% ------------------------------------------------------------
\medskip
\noindent (i) \textit{Observations on the wave numbers of the disturbance manifested at
marginal stability}

In some ways, the most striking phenomenon observed in these
experiments is the sudden appearance of regularly spaced toroidal
vortices of rectangular cross-section.  An example of such manifestation,
as first photographed by Taylor, is shown in Fig.~82.

The height of the vortices, which appear at marginal stability, is
related to the wave number $a$ introduced in the theory by
\begin{equation}\label{eq:7-335}
\text{height of vortex} = R_2\,\pi/a.
\end{equation}
From the measured heights, the wave number $a$ can be deduced.  In
Fig.~81 a comparison is made between the experimentally deduced and
the theoretically predicted values of~$a$.

\figurePlaceholder{81}{Comparison of the calculated (full-line curve and the dots with
vertical lines) and observed (dots) wave numbers at the onset of instability.}

Photographs of the vortices obtained by Donnelly and Fultz under
various conditions are shown in Figs.~83--84.

The onset of instability at $\mu = 0$ is illustrated by the sequence of
photographs $a$, $b$, and $c$ in Fig.~82: ($a$)~shows the liquid in laminar flow
at a speed slightly below critical; ($b$)~shows the beginnings of three-dimensional
motion when the ink gathering together starts to move

\figurePlaceholder{85}{Comparison between the observed and the
predicted $(\Omega_1, \Omega_2)$-relation for the onset of instability
for the case $\eta = \frac{1}{2}$: the dots represent the points at
which instability was observed to occur in the experi\-ments
of Donnelly and Fultz and the full-line curve
represents the theoretical relation.}

\noindent radially outwards to form the first cell; and ($c$)~shows the vortex fully
developed from the initial state shown in ($b$) with no additional change
in the speed.

When both cylinders are rotated in the same direction, the appearance
of the cells is much the same as in the case $\mu = 0$.  An example is shown
in Fig.~82\,$d$.

When the cylinders are rotated in opposite directions, the phenomena
one observes are much more complex.  Theoretically, we expect two cells
in each compartment, with the motion in the outer cell relatively much
weaker (see Fig.~71\,$a$, $b$).  Since the ink is ejected from holes in the inner
cylinder, we should not, under the circumstances, expect to witness
motions in the outer cell.  Also, for reasons which we have explained in
\S\S\,68 and 71\,($b$), we must expect that, as $\mu \to -\infty$, the cells at marginal
stability appear very close to the inner cylinder and become increasingly
tiny.  All these expectations are strikingly confirmed by the sequence of
photographs in Figs.~83 and 84.

\figurePlaceholder{82}{The onset of instability with the outer cylinder at rest, $\mu = 0$. $P_c = 4{\cdot}493$~sec;
(a) laminar flow, $P = 4{\cdot}500$~sec; (b) beginning of radial motion at $P = 4{\cdot}483$~sec;
(c) Appearance of cells with outer cylinder at rest: cells at marginal stability, $P = P_c = 4{\cdot}466$~sec.
(d) Appearance of cells with cylinders rotating in the same direction: $P = 3{\cdot}544$~sec, $\mu = 0{\cdot}1184$.}

\figurePlaceholder{83}{Photographs of cells with cylinders ($\eta = \frac{1}{2}$) rotating in opposite directions:
(a) $\mu = -0{\cdot}4116$; (b) $\mu = -0{\cdot}5$; and (c) $\mu = -1{\cdot}74$.}

\figurePlaceholder{84}{Photographs of cells with cylinders ($\eta = \frac{1}{2}$) rotating in opposite directions:
(a) $\mu = -2{\cdot}44$; (b) $\mu = -3{\cdot}02$; (c) $\mu = -5{\cdot}86$; and (d) $\mu = -6{\cdot}83$.}

%% ------------------------------------------------------------
%% (ii) Comparison with predicted relations
%% ------------------------------------------------------------
\medskip
\noindent (ii) \textit{Comparison between the measured and the predicted $(T_c, \mu)$-relations}

The results of Donnelly and Fultz, obtained with their apparatus on
the critical angular velocities at which instability occurs, are exhibited
in Figs.~85 and 86 as plots analogous to those of Figs.~72 and 74.  The
agreement of the experimental results with the calculations of \S\,73\,($b$)
is very satisfactory.

\figurePlaceholder{86}{The variation of the critical Taylor number $(T_c)$ for the
case $\eta = \frac{1}{2}$ as a function of $\kappa\;(= (1-4\mu)/(1-\mu))$.  The solid line repre\-sents
the theoretically calculated relation and the dashed line is a
linear extrapolation of it.  The dots represent the points at which
instability was observed to occur.  Data to the left of the arrow are
unduly influenced by small errors in $\mu$.}

As we have already remarked, one of Taylor's original sets of measurements
was obtained by an apparatus in which the gap width was a small
fraction ($\sim\!5$ per cent) of the radius of the outer cylinder.  These measurements
can, therefore, be compared with the theoretical results of \S\,71\,($b$).
This comparison is included in Fig.~87 (curve labelled $\eta = 0{\cdot}8418$).
The agreement of the theoretical curve with the results of the measurements
is good over the entire range.

Fig.~87 includes measurements by Taylor for other values of $\eta$; and
also of Donnelly and Fultz for values of $-\mu$ beyond the range of the
calculations of \S\,73\,($b$).  From an examination of the data presented in
Fig.~87, Donnelly and Fultz have inferred that an asymptotic relation
of the form,
\begin{equation}\label{eq:7-336}
(\Omega_1\,R_1^2/\nu)^5 \to \text{constant}\times (|\Omega_2|\,R_2^2/\nu)^3
\quad\text{for}\quad \Omega_2/\Omega_1 \to -\infty,
\end{equation}
exists.

\figurePlaceholder{87}{Comparisons between the observed and the predicted $(\omega_1, \omega_2)$-relations for
the onset of stability: $\omega_1 = \Omega_1\,R_1^2/\nu$ and $\omega_2 = \Omega_2\,R_2^2/\nu$.  The uppermost solid curve
and the one at the extreme left represent the theoretical relations derived in \S\,71\,($b$)
and \S\,73\,($b$), respectively; the dashed-line continuation of the latter represents an
extrapolation of the theoretical relation.  The remaining solid lines are derived from
the asymptotic relation~\eqref{eq:7-342}.  The experimental data for $\eta = 0{\cdot}9418$, $0{\cdot}8798$, and
$0{\cdot}7435$ are those of G.~I.~Taylor; the data for $\eta = 0{\cdot}5$ are those of Donnelly and Fultz.}

Donnelly and Fultz have pointed out that the asymptotic relation
\eqref{eq:7-336} has a simple physical interpretation.  Arguing as in \S\,71\,(e), they
first observe that the onset of instability, in the case when $\mu < 0$, is
primarily to be associated with the violation of Rayleigh's criterion in
the layers immediately adjoining the inner cylinder and extending not
much beyond the nodal surface (at $r = R_n$, say).  Since the nodal
surface approaches the inner cylinder as $\mu \to -\infty$, the fluid which
becomes unstable is essentially confined to a narrow gap between $R_1$
and $R_n$.  Accordingly, we may write (cf.\ equation~(325))
\begin{equation}\label{eq:7-337}
\frac{4\Omega_1^2}{\nu^2}\,\frac{R_1^2\,d_0^4}{R_n^2 - R_1^2} = T_0
\end{equation}
as the condition for the onset of instability, where $d_0 = R_n - R_1$ and
$T_0$ is a number of the order of 1000.  Since the present arguments pre\-suppose
that $d_0 \ll R_1$, an equivalent form of~\eqref{eq:7-337} is
\begin{equation}\label{eq:7-338}
\left(\frac{\Omega_1\,R_1^2}{\nu}\right)^{\!2} = \tfrac{1}{2}\,T_0\!\left(\frac{R_1}{d_0}\right)^{\!3}\!.
\end{equation}
But, according to equation~(29),
\begin{equation}\label{eq:7-339}
d_0 = R_2\,\eta\!\left(\frac{1+|\mu|}{\eta^2+|\mu|}\right)^{\!1/2} - R_1
= R_1\!\left(\sqrt{\frac{1+|\mu|}{\eta^2+|\mu|}} - 1\right)\!.
\end{equation}
For $\mu \to -\infty$,
\begin{equation}\label{eq:7-340}
\frac{d_0}{R_1} \to \frac{1-\eta^2}{2|\mu|}
= \frac{1}{2}\left|\frac{\Omega_1}{\Omega_2}\right|(1-\eta^2).
\end{equation}
Inserting this value of $d_0/R_1$ in equation~\eqref{eq:7-338}, we obtain
\begin{equation}\label{eq:7-341}
\left(\frac{\Omega_1\,R_1^2}{\nu}\right)^{\!2}
= \frac{4\,T_0}{(1-\eta^2)^3}\left|\frac{\Omega_2}{\Omega_1}\right|^{\!3}
\qquad (\mu \to -\infty),
\end{equation}
or, alternatively,
\begin{equation}\label{eq:7-342}
\left(\frac{\Omega_1\,R_1^2}{\nu}\right)^{\!5}
= \frac{4\,\eta^6\,T_0}{(1-\eta^2)^3}\left(\frac{|\Omega_2|\,R_2^2}{\nu}\right)^{\!3}
\qquad (\mu \to -\infty),
\end{equation}
in agreement with~\eqref{eq:7-336}.

The values of $T_0$, empirically determined by Donnelly and Fultz for
various values of $\eta$, are
\begin{equation}\label{eq:7-343}
T_0 = \begin{cases}
450  & \eta = 0{\cdot}9418 \\
530  & \eta = 0{\cdot}8798 \\
670  & \eta = 0{\cdot}7435 \\
1{,}180 & \eta = 0{\cdot}5.
\end{cases}
\end{equation}

%% ============================================================
%% BIBLIOGRAPHICAL NOTES
%% ============================================================
\section*{BIBLIOGRAPHICAL NOTES}

The fundamental papers on the stability of Couette flow are those of Lord
Rayleigh and Sir Geoffrey Taylor:

\begin{enumerate}
\item \textsc{Lord Rayleigh}, `On the dynamics of revolving fluids', \textit{Scientific Papers},
\textbf{6}, 447--53, Cambridge, England, 1920.

\item G.~I. \textsc{Taylor}, `Stability of a viscous liquid contained between two rotating
cylinders', \textit{Phil.\ Trans.\ Roy.\ Soc.\ (London)} A, \textbf{223}, 289--343 (1923).

\noindent In reference 1, Rayleigh treated the case of inviscid flow and derived the general
criterion known after him.  In reference~2, Taylor treated for the first time the
corresponding problem in viscid flow, both theoretically and experimentally.

\noindent The following general reference, in which the problem of hydrodynamic stability
is broadly surveyed with special reference to Couette flow, may be noted:

\item J.~L. \textsc{Synge}, `Hydrodynamic stability', 227--69, \textit{Semicentennial Addresses
of the American Mathematical Society}, ii, New York, 1938.

\item C.~C. \textsc{Lin}, \textit{The Theory of Hydrodynamic Stability}, Cambridge, England, 1955.

\noindent We may also refer to the relevant sections in:

\medskip
\noindent\S\,66.  As we have already remarked, Rayleigh developed his criterion for the
stability of inviscid rotational flow in reference~1.  In his discussion, Rayleigh
argued largely in terms of the analogy between the present problem and the
problem of the stability of an incompressible fluid of variable density.  For his
discussion of the latter problem, see:

\item Lord \textsc{Rayleigh}, `Investigation of the character of the equilibrium of an
incompressible heavy fluid of variable density', \textit{Scientific Papers}, \textbf{2},
200--7, Cambridge, England, 1900.

\medskip
\noindent\S\,67.  The analytical derivation of Rayleigh's criterion (for the case $m = 0$)
is due to:

\item J.~L. \textsc{Synge}, `The stability of heterogeneous liquids', \textit{Trans.\ of the Royal
Society of Canada}, \textbf{27}, 1--18 (1933).

\noindent The general case ($m \neq 0$) is considered in:

\item S. \textsc{Chandrasekhar}, `The stability of inviscid flow between rotating
cylinders', \textit{J.\ Indian Math.\ Soc.} (in press).

\noindent The discussion in this section is largely derived from reference~7.

\noindent A useful reference for the Sturmian theory and its later development is:

\item E.~L. \textsc{Ince}, \textit{Ordinary Differential Equations}, chapter~x, Longmans, Green \&
Co.\ Ltd., London, 1927.

\medskip
\noindent\S\,68.  The oscillations of a rotating column of liquid were first considered by
Lord Kelvin:

\item Lord \textsc{Kelvin}, `Vibrations of a columnar vortex', 152--65, \textit{Mathematical
and Physical Papers}, iv, \textit{Hydrodynamics and General Dynamics}, Cambridge,
England, 1910.

\noindent An extensive treatment of this subject with meteorological overtones will be
found in:

\item V. \textsc{Bjerknes}, J. \textsc{Bjerknes}, H. \textsc{Solberg}, and T. \textsc{Bergeron}, \textit{Physikalische
Hydrodynamik}, chapter~11, Springer, Berlin, 1933.

\medskip
\noindent\S\,69\,(a).  Fultz first showed how the natural modes of oscillation of a rotating
cylinder of liquid can be excited and maintained.

\item D. \textsc{Fultz}, `A note on overstability and the elasto-inertia oscillations of
Kelvin, Solberg, and Bjerknes', \textit{J.\ Meteorology}, \textbf{16}, 199--208 (1959).

\noindent The solution of equation~(96) for different values of $m$ was carried out by Miss
Donna Elbert; the results of her calculations are included in Fig.~63.
The roots listed in Table~XXXI (with the exception of those for $\eta = 0$) are taken
from:

\item S. \textsc{Chandrasekhar} and Donna \textsc{Elbert}, `The roots of
\[
Y_n(\lambda\eta)J_n(\lambda) - J_n(\lambda\eta)Y_n(\lambda) = 0\text{'},
\]
\textit{Proc.\ Camb.\ Phil.\ Soc.} \textbf{50}, 266--8 (1954).

\medskip
\noindent\S\,68\,(b).  The solution for the narrow gap described in this section is due to Reid:

\item W.~H. \textsc{Reid}, `Inviscid modes of instability in Couette flow', \textit{J.\ Math.\
Analysis and Applications}, \textbf{1}, 411--22 (1960).

\noindent I am grateful to Dr.~Reid for providing me with copies for the illustrations
(Figs.~65--67) included in this section.

\noindent A convenient reference for Airy's functions, in terms of which the solution for
this case is found, is:

\item J.~C.~P. \textsc{Miller}, `The Airy integral, giving tables of solutions of the
differential equation $y'' = xy$', \textit{British Assoc.\ Math.\ Tables}, Part-volume~B,
Cambridge, England, 1946.

\medskip
\noindent\S\,69.  The illustration (Fig.~68) in this section is taken from Taylor's original
paper (reference~2).

\noindent\S\,70.  The theorem proved in \S\,(a) is due to Synge.

\item J.~L. \textsc{Synge}, `On the stability of a viscous liquid between rotating coaxial
cylinders', \textit{Proc.\ Roy.\ Soc.\ (London)} A, \textbf{167}, 250--6 (1938).

\medskip
\noindent\S\,71\,(a).  The analysis in this section is based on:

\item S. \textsc{Chandrasekhar}, `The stability of viscous flow between rotating cylin\-ders',
\textit{Mathematika}, \textbf{1}, 5--13 (1954).

\noindent Alternative treatments of the characteristic value problem presented by equa\-tions~(200)--(202) have been given by:

\item D. \textsc{Meksyn}, `Stability of viscous flow between rotating cylinders.~I',
\textit{Proc.\ Roy.\ Soc.\ (London)} A, \textbf{187}, 115--28 (1946).

\item \leavevmode\kern-0.3em---\kern0.3em `Stability of viscous flow between rotating cylinders.  II. Cylinders
rotating in opposite directions', ibid.\ 480--91 (1946).

\item \leavevmode\kern-0.3em---\kern0.3em `Stability of viscous flow between rotating cylinders.  III.~Integra\-tion
of a sixth order linear equation', ibid.\ 492--504 (1946).

\item R.~C. \textsc{Di Prima}, `Application of the Galerkin method to problems in
hydrodynamic stability', \textit{Quart.\ Appl.\ Math.} \textbf{13}, 55--62 (1955).

\noindent The methods used by these authors (as well as by Taylor in reference~2) do not
enable as systematic (or as accurate) a solution of the problem as the method
described in the text.

\noindent An extension of the analysis in reference~16 to include a term in $\varepsilon^2$ on the right-hand
side of equation~(201) has been given by:

\item H. \textsc{Steinman}, `The stability of viscous flow between rotating cylinders',
\textit{Quart.\ Appl.\ Math.} \textbf{14}, 27--33 (1956).

\medskip
\noindent\S\,71\,(b).  The numerical results given in reference~16 have been supplemented.
These additional results (even as the original ones in reference~16) are due to
Miss Donna Elbert.

\medskip
\noindent\S\,71\,(c).  The alternative method described in this section has not been published
before.

\medskip
\noindent\S\,71\,(d).  The relation disclosed in this section between the present problem and
the simple B\'enard problem was first noticed by Low:

\item A.~R. \textsc{Low}, `Instability of viscous fluid motion', \textit{Nature}, \textbf{115}, 299--300 (1925).

\noindent The analytical basis for this identity in the characteristic values of the two
problems (in a certain approximation) was provided by:

\item H. \textsc{Jeffreys}, `Some cases of instability in fluid motion', \textit{Proc.\ Roy.\ Soc.\
(London)} A, \textbf{118}, 195--208 (1928).

\noindent The perturbation procedure described in this section is, however, new.  I am
indebted to Dr.~F.~Vandervort for the evaluation of the matrix elements in the
solution~(252).

\medskip
\noindent\S\,71\,(e).  The essential arguments of this section can be found (in varying degrees
of explicitness) in the writings of Meksyn (reference~18), Lin (reference~4), Di
Prima (reference~20), and Donnelly and Fultz (reference~28 below).

\medskip
\noindent\S\,72.  The discussion of overstability in this section has not been published before.

\medskip
\noindent\S\,73.  The analysis in this section is based on:

\item S. \textsc{Chandrasekhar}, `The stability of viscous flow between rotating cylin\-ders',
\textit{Proc.\ Roy.\ Soc.\ (London)} A, \textbf{246}, 301--11 (1958).

\noindent The method of solution derives from:

\item S. \textsc{Chandrasekhar} and W.~H. \textsc{Reid}, `On the expansion of functions which
satisfy four boundary conditions', \textit{Proc.\ Nat.\ Acad.\ Sci.} \textbf{43}, 521--7 (1957).

\noindent The orthogonal functions, $\Psi_i(a, r)$, in terms of which the solution is found, are
tabulated in:

\item S. \textsc{Chandrasekhar} and Donna~D. \textsc{Elbert}, `On orthogonal functions
which satisfy four boundary conditions.  III.~Tables for use in Fourier--Bessel-type
expansions', \textit{Astrophys.\ J.\ Supp.\ Ser.} \textbf{3}, 453--8 (1958).

\medskip
\noindent\S\,74.  The references to the experimental work described in this section are:

\item R.~J. \textsc{Donnelly}, `Experiments on the stability of viscous flow between
rotating cylinders.  I.~Torque measurements', \textit{Proc.\ Roy.\ Soc.\ (London)} A,
\textbf{246}, 312--25 (1958).

\item \leavevmode\kern-0.3em---\kern0.3em and D. \textsc{Fultz}, `Experiments on the stability of viscous flow between
rotating cylinders.  II.~Visual observations', ibid.\ \textbf{258}, 101--23 (1960).

\noindent The results of similar and related experiments will be found in:

\item J.~W. \textsc{Lewis}, `An experimental study of the motion of a viscous liquid
contained between two coaxial cylinders', \textit{Proc.\ Roy.\ Soc.\ (London)} A, \textbf{117},
388--407 (1928).

\item T. \textsc{Terada} and T. \textsc{Hattori}, `Some experiments on motions of fluids, IV',
\textit{Rep.\ Aero.\ Res.\ Inst.\ Tokyo}, \textbf{2}, 287--326 (1926).

\item F. \textsc{Wendt}, `Turbulente Str\"omungen zwischen zwei rotierenden koaxialen
Zylindern', \textit{Ingen.\ Arch.} \textbf{4}, 577--96 (1933).

\item G.~I. \textsc{Taylor}, `Fluid friction between rotating cylinders.  I.~Torque
measurements', \textit{Proc.\ Roy.\ Soc.\ (London)} A, \textbf{157}, 546--64 (1936).

\item \leavevmode\kern-0.3em---\kern0.3em `Fluid friction between rotating cylinders. II. Distribution of
velocity between concentric cylinders when outer one is rotating and inner
one is at rest', ibid.\ 565--78 (1936).

\item F. \textsc{Schultz-Grunow} and H. \textsc{Hein}, `Beitrag zur Couettestr\"omung', \textit{Z.\ f.\
Flugwiss.} \textbf{4}, 28--30 (1956).

\end{enumerate}
